\documentclass[11pt,a4paper]{article}
\usepackage[utf8]{inputenc}
\usepackage[T1]{fontenc}
\usepackage{lmodern}
\usepackage{geometry}
\usepackage{amsmath,amssymb,amsthm}
\usepackage{booktabs}
\usepackage[hidelinks]{hyperref}
\usepackage{url}
\geometry{margin=2.5cm}
\emergencystretch=3em

\newtheorem{lemma}{Lemma}
\newtheorem{target}{Analytic Target}
\newtheorem{observation}{Observation}

\title{\bfseries Derived Prime-Harmonic Envelope on CA Support\\
\large Supplement to the Corrected Robin--MVDC Programme}
\author{Ing.~Robert~Polak \\ \small robopol@gmail.com \\ \small \url{https://robopol.sk}}
\date{\today}

\begin{document}
\maketitle

\begin{abstract}
This supplement continues the corrected status PDF manuscript of the
Robin--MVDC programme \cite{polakcorrected2026}. It records the exact point
reached by the reciprocal-prime Route II. First, it recalls the certified
\(B_n^{\rm env}\)-envelope for \(\sigma(n)/n\), explains how it is obtained
from the exact divisor product, and shows how that envelope produces the
Robin reserve used by the block ledger. Starting from the first-moment beta
ledger on CA support blocks, the remaining condition is then rewritten as an
upper envelope for the prime-harmonic remainder
\[
  A(x)-B
  =
  \sum_{p\le x}\frac1p-\log\log x-B.
\]
The resulting admissible envelope has the form
\[
  A(x_i)-B
  \le
  \frac{C_{\rm req}(x_i)}{\sqrt{x_i}\log x_i}
\]
on CA support endpoints \(x_i\), where \(C_{\rm req}(x_i)\) is determined by
the Robin--MVDC ledger. This is a derived requirement, not a proved theorem
about the primes. The remaining analytic obstruction is therefore the
infinite-range certification of this prime-harmonic oscillation bound on the
CA support sequence. Existing explicit reciprocal-prime estimates are reviewed
only to clarify why the known all-range bounds are too coarse for this
ledger.
\end{abstract}

\paragraph{Keywords:} Riemann Hypothesis; Robin's inequality; colossally
abundant numbers; beta envelope; reciprocal primes; prime-harmonic remainder;
MVDC; explicit prime estimates.

\noindent\emph{Status of this supplement. This is a working research note, not a
completed proof. It depends on the exact Robin ledger of the corrected status
PDF manuscript \cite{polakcorrected2026}. It records the exact point reached by
the Robin--MVDC Route II reduction: the remaining Robin ledger is reduced to a
certified upper envelope for the prime-harmonic fluctuation on CA support
endpoints. The envelope is derived by the ledger algebra; the missing theorem
is the infinite-range prime-oscillation bound needed to certify that the primes
obey it.}

\section{Relation To The Main PDF}

This note is a supplement to the corrected status PDF manuscript
\cite{polakcorrected2026}, specifically to its exact Robin ledger and to
Route II, the first-moment reciprocal-prime certificate.
The present text does not replace that manuscript. Its role is narrower: it
spells out the \(B_n\)-envelope used to dominate \(\sigma(n)/n\), then follows
the Route II algebra until the remaining prime-harmonic oscillation target is
made explicit.

The main PDF denotes the actual multiplicative divisor deficit by \(A(n)\).
In this supplement we write that deficit as \(\mathcal D(n,x)\), because the
letter \(A(x)\) is reserved below for the reciprocal-prime remainder
\[
  A(x)=\sum_{p\le x}\frac1p-\log\log x.
\]
We also reserve \(B\) for the Meissel--Mertens prime constant. The divisor-sum
upper envelope is therefore written as \(B_n^{\rm env}\).

\section{The \(B_n\) Upper Envelope}

Robin's criterion \cite{robin1984} reduces the Riemann Hypothesis to Robin's
inequality above the finite cutoff \(5040\).  The present programme tests the
remaining large-candidate ledger on CA/SA-type support profiles, in the spirit
of the extremal divisor-sum structure studied by Alaoglu and Erdos
\cite{alaogluerdos1944}. For the CA/SA profiles retained in the corrected
ledger, the support is full up to the largest prime \(x=P(n)\):
\[
  n=\prod_{p\le x}p^{a_p},\qquad a_p\ge1.
\]
For each prime power,
\[
  \frac{\sigma(p^a)}{p^a}
  =
  1+\frac1p+\cdots+\frac1{p^a}
  =
  \frac{1-p^{-(a+1)}}{1-p^{-1}}.
\]
Hence
\[
  \frac{\sigma(n)}{n}
  =
  \prod_{p\le x}
  \frac{1-p^{-(a_p+1)}}{1-p^{-1}}
  =
  \beta(x)\exp\{-\mathcal D(n,x)\},
\]
where
\[
  \beta(x):=\prod_{p\le x}\frac{p}{p-1},
  \qquad
  \mathcal D(n,x):=
  -\sum_{p\le x}\log\left(1-p^{-(a_p+1)}\right).
\]

Let \(\mathcal D_{\rm env}(n,x)\) be any certified lower envelope for the
actual deficit:
\[
  \mathcal D_{\rm env}(n,x)\le \mathcal D(n,x).
\]
Define the divisor-sum envelope
\begin{equation}\label{eq:bn-env}
  B_n^{\rm env}:=
  \beta(x)\exp\{-\mathcal D_{\rm env}(n,x)\}.
\end{equation}
This is the \(B_n\)-envelope used by the ledger.

\begin{lemma}[Certified divisor-sum envelope]\label{lem:bn-envelope}
If \(\mathcal D_{\rm env}(n,x)\le \mathcal D(n,x)\), then
\[
  \frac{\sigma(n)}{n}\le B_n^{\rm env}.
\]
\end{lemma}

\begin{proof}
The map \(t\mapsto e^{-t}\) is decreasing. Therefore
\[
  e^{-\mathcal D(n,x)}
  \le
  e^{-\mathcal D_{\rm env}(n,x)}.
\]
Multiplying by \(\beta(x)>0\) and using the exact product formula above gives
\[
  \frac{\sigma(n)}{n}
  =
  \beta(x)e^{-\mathcal D(n,x)}
  \le
  \beta(x)e^{-\mathcal D_{\rm env}(n,x)}
  =
  B_n^{\rm env}.
\]
\end{proof}

Thus \(B_n^{\rm env}\) is not a fitted numerical curve. It is obtained by
starting from the exact formula for \(\sigma(n)/n\) and replacing the actual
divisor deficit by a certified lower bound. The trivial envelope is
\(\beta(x)\), obtained from \(\mathcal D_{\rm env}=0\). Any stronger certified
lower envelope for \(\mathcal D(n,x)\) lowers \(B_n^{\rm env}\) while keeping
it above \(\sigma(n)/n\).

Robin's inequality follows from the sufficient envelope condition
\[
  B_n^{\rm env}\le e^\gamma\log\log n.
\]
Taking logarithms and writing
\[
  E(x):=\log\beta(x)-\gamma-\log\log x,
  \qquad
  B_{\log}(n,x):=
  \log\frac{\log\log n}{\log x},
\]
we obtain
\[
  \log B_n^{\rm env}
  -\gamma-\log(\log\log n)
  =
  E(x)-\mathcal D_{\rm env}(n,x)-B_{\log}(n,x).
\]
Therefore it is sufficient to prove the exact reserve gate
\begin{equation}\label{eq:reserve-gate}
  E(x)\le
  \mathcal D_{\rm env}(n,x)+B_{\log}(n,x).
\end{equation}
The reserve used below is precisely a certified lower bound for the right-hand
side of \eqref{eq:reserve-gate} on the CA endpoint under consideration.

\section{Block Objects}

Let \(x_i\) denote the sampled CA support endpoints and
write
\[
  Y=x_{i-1},\qquad x=x_i.
\]
The reciprocal-prime remainder is
\[
  A(x):=\sum_{p\le x}\frac1p-\log\log x,
\]
and \(B\) denotes the Meissel--Mertens prime constant. The fluctuation to be
controlled is therefore
\[
  A(x)-B.
\]

On a CA support block put
\[
  \nu=\pi(x)-\pi(Y),\qquad
  H=\log\frac{\log x}{\log Y},\qquad
  \mu=\frac{H}{\nu}.
\]
Let \(R(x)\) be the certified CA reserve at the endpoint \(x\), and let
\(U(Y)\) be the incoming certified upper ledger at \(Y\). The available block
threshold is
\[
  T(Y,x):=R(x)-U(Y).
\]

\section{First-Moment Gate}

The beta-error block is
\[
  Q(Y,x)
  =
  \sum_{Y<p\le x}-\log(1-1/p)
  -
  \log\frac{\log x}{\log Y}.
\]
With
\[
  u_p=e^{-\mu}\frac{p}{p-1}-1,
\]
we have
\[
  Q(Y,x)=\sum_{Y<p\le x}\log(1+u_p)\le \sum_{Y<p\le x}u_p=:M_1(Y,x).
\]
Writing
\[
  S_{-1}(Y,x):=\sum_{Y<p\le x}\frac1{p-1},
\]
the first moment is exactly
\begin{equation}\label{eq:m1-exact}
  M_1(Y,x)=e^{-\mu}\bigl(\nu+S_{-1}(Y,x)\bigr)-\nu.
\end{equation}

Thus Route II is closed on the block if the sufficient first-moment gate
\begin{equation}\label{eq:m1-gate}
  M_1(Y,x)\le T(Y,x)
\end{equation}
holds. Solving \eqref{eq:m1-gate} with \eqref{eq:m1-exact} gives
\begin{equation}\label{eq:sminus-gate}
  S_{-1}(Y,x)
  \le
  \nu(e^\mu-1)+e^\mu T(Y,x).
\end{equation}

\section{Reduction To \(A(x)\)}

The key exact identity is
\[
  \frac1{p-1}=\frac1p+\frac1{p(p-1)}.
\]
Therefore
\begin{align*}
  S_{-1}(Y,x)
  &=
  \sum_{Y<p\le x}\frac1p
  +
  \sum_{Y<p\le x}\frac1{p(p-1)}\\
  &=
  \log\frac{\log x}{\log Y}+A(x)-A(Y)+C_2(Y,x),
\end{align*}
where
\[
  C_2(Y,x):=\sum_{Y<p\le x}\frac1{p(p-1)}.
\]

Substituting this into \eqref{eq:sminus-gate} gives
\[
  H+A(x)-A(Y)+C_2(Y,x)
  \le
  \nu(e^\mu-1)+e^\mu T(Y,x).
\]
Solving for \(A(x)\), the first-moment gate is equivalent to
\begin{equation}\label{eq:a-req-direct}
  A(x)\le A_{\rm req}(Y,x),
\end{equation}
where
\begin{equation}\label{eq:a-req}
  A_{\rm req}(Y,x)
  :=
  A(Y)-C_2(Y,x)+\nu(e^\mu-1)-H+e^\mu T(Y,x).
\end{equation}

Equivalently, define
\[
  D(Y,x):=\nu\left(1-(1+\mu)e^{-\mu}\right).
\]
Since
\[
  e^\mu D(Y,x)=\nu(e^\mu-1)-H,
\]
the same required envelope can be written as
\begin{equation}\label{eq:a-req-d}
  A_{\rm req}(Y,x)
  =
  A(Y)-C_2(Y,x)+e^\mu\bigl(T(Y,x)+D(Y,x)\bigr).
\end{equation}

\section{The Derived Envelope}

Define the normalized required constant by
\begin{equation}\label{eq:c-req}
  C_{\rm req}(x)
  :=
  \bigl(A_{\rm req}(Y,x)-B\bigr)\sqrt{x}\log x.
\end{equation}
Then \eqref{eq:a-req-direct} is exactly the same statement as
\begin{equation}\label{eq:derived-envelope}
  A(x)-B
  \le
  \frac{C_{\rm req}(x)}{\sqrt{x}\log x}.
\end{equation}

This is the point reached by the reduction. The right-hand side in
\eqref{eq:derived-envelope} is not inserted by hand; it is the admissible
fluctuation budget produced by the Robin--MVDC ledger. It says how small the
prime-harmonic fluctuation must be on the CA support endpoint \(x\) for the
Robin ledger to close.

What has not been proved is that the primes satisfy
\eqref{eq:derived-envelope} for all sufficiently large CA support endpoints.
Thus the remaining obstruction is the prime-core problem:
\[
  \text{certify the positive oscillation of }A(x)-B
  \text{ on the infinite CA support range.}
\]

\section{Why A Fixed \(C\)-Scale Matters}

The hard content is the scale
\[
  \frac{1}{\sqrt{x}\log x},
\]
not the particular numerical value of a constant. A theorem of the form
\[
  A(x_i)-B\le \frac{C_*}{\sqrt{x_i}\log x_i}
\]
with a fixed constant \(C_*\) on all sufficiently large CA support endpoints
would be a genuine RH-scale oscillation control. If, in addition,
\[
  C_*\le C_{\rm req}(x_i)
\]
eventually, or if the excess \(C_*-C_{\rm req}(x_i)\) is absorbed by a
cumulative ledger margin, then Route II would close.

By contrast, a bound that can be rewritten only as
\[
  A(x)-B\le \frac{C_{\rm eff}(x)}{\sqrt{x}\log x}
\]
with \(C_{\rm eff}(x)\to\infty\) does not provide the needed oscillation
control. It merely expresses a weaker estimate in the same notation.

\section{Literature Check}

Rosser--Schoenfeld \cite{rosser1962} prove the sharp finite-window
estimate
\[
  A(x)\le B+\frac{2}{\sqrt{x}\log x}
\]
for \(1<x<10^8\). This is the origin of the useful \(C=2\) benchmark, but it is
not an infinite-range input.

Their all-large reciprocal-prime estimate is of logarithmic scale,
\[
  |A(x)-B|\le \frac{1}{2\log^2 x},
\]
which corresponds in the present normalization to an effective constant
\[
  C_{\rm eff}(x)=\frac{\sqrt{x}}{2\log x}.
\]
This grows with \(x\), and is therefore too coarse for the Robin--MVDC ledger.

Dusart \cite{dusart2018} and Axler \cite{axler2018} give sharper explicit
reciprocal-prime estimates in powers of \(\log x\). For example Axler's upper
bound
\[
  A(x)-B
  \le
  \frac{1}{20\log^3 x}
  +
  \frac{3}{16\log^4 x}
  \qquad (x\ge 46\,909\,074)
\]
becomes
\[
  C_{\rm eff}(x)
  =
  \sqrt{x}
  \left(
    \frac{1}{20\log^2 x}
    +
    \frac{3}{16\log^3 x}
  \right)
\]
in the \(C/(\sqrt{x}\log x)\) scale. This is useful as a finite diagnostic, but
it is not a fixed-\(C\) RH-scale theorem.

\section{Numerical Audit}

The support script
\[
  \texttt{current\_support/robin\_mvdc\_status\_support/}
  \texttt{test\_c2\_shortfall\_trend.py}
\]
computes \(C_{\rm req}\) and the measured value
\[
  C_{\rm actual}(x):=\bigl(A(x)-B\bigr)\sqrt{x}\log x
\]
on sampled CA support endpoints. The following table is numerical evidence
only; it is not a proof of the infinite-range oscillation bound. The
corresponding auxiliary scripts and audit data are maintained in the support
repository \cite{polakriemannrepo2026}.

\[
\begin{array}{r|r|r|r|r}
x & C_{\rm req}(x) & C_{\rm actual}(x) &
C_{\rm req}(x)-C_{\rm actual}(x) & C_{\rm req}(x)-2\\
\hline
6382007   & 1.563181 & 0.770993 & 0.792188 & -0.436819\\
12253883  & 1.768796 & 0.976166 & 0.792630 & -0.231204\\
29093377  & 1.629640 & 0.836043 & 0.793597 & -0.370360\\
56048351  & 1.477260 & 0.697665 & 0.779596 & -0.522740\\
78580847  & 1.657538 & 0.886993 & 0.770545 & -0.342462\\
108115627 & 1.789141 & 0.996545 & 0.792596 & -0.210859\\
175589599 & 2.044976 & 1.253114 & 0.791862 & +0.044976\\
258144067 & 1.856973 & 1.071509 & 0.785464 & -0.143027\\
339396977 & 1.438563 & 0.663823 & 0.774740 & -0.561437\\
499283831 & 2.029917 & 1.272876 & 0.757041 & +0.029917\\
701361907 & 1.917137 & 1.145477 & 0.771660 & -0.082863\\
966679613 & 1.699811 & 0.910342 & 0.789469 & -0.300189
\end{array}
\]

The measured fluctuation is well below the derived admissible envelope on the
sampled range. The difficulty is not the finite arithmetic in this table, but
the missing proof that this remains true on the infinite CA support range.

\section{Final Analytic Target}

\begin{target}[CA-support prime-harmonic envelope]
Prove that for all sufficiently large CA support endpoints \(x_i\),
\[
  A(x_i)-B
  \le
  \frac{C_{\rm req}(x_i)}{\sqrt{x_i}\log x_i},
\]
where \(C_{\rm req}(x_i)\) is the admissible ledger constant defined by
\eqref{eq:c-req}.
\end{target}

This target is exactly what remains after the Route II reduction. It is not an
MVDC algebra problem anymore; it is a certified oscillation-control problem for
the reciprocal-prime remainder on the CA support sequence.

\begin{thebibliography}{99}

\bibitem{polakcorrected2026}
R.~Polak,
\newblock Corrected Status of the Robin--MVDC Programme:
Superseding the Previous Robin--Mertens Proof Branch and Consolidated Manuscript,
\newblock PDF working manuscript, 2026.

\bibitem{robin1984}
G.~Robin,
\newblock Grandes valeurs de la fonction somme des diviseurs et hypothese de Riemann,
\newblock \emph{Journal de Mathematiques Pures et Appliquees}, 63 (1984), 187--213.

\bibitem{alaogluerdos1944}
L.~Alaoglu and P.~Erdos,
\newblock On highly composite and similar numbers,
\newblock \emph{Transactions of the American Mathematical Society}, 56 (1944), 448--469.

\bibitem{rosser1962}
J.~B.~Rosser and L.~Schoenfeld,
\newblock Approximate formulas for some functions of prime numbers,
\newblock \emph{Illinois Journal of Mathematics}, 6 (1962), 64--94.

\bibitem{dusart2018}
P.~Dusart,
\newblock Explicit estimates of some functions over primes,
\newblock \emph{Ramanujan Journal}, 45 (2018), 227--251.

\bibitem{axler2018}
C.~Axler,
\newblock New estimates for some functions defined over primes,
\newblock Integers 18 (2018), Paper No. A52, 21 pp.

\bibitem{polakriemannrepo2026}
R.~Polak,
\newblock Riemann-hypothesis: auxiliary scripts and numerical audit data for
the Robin--MVDC programme,
\newblock GitHub repository, 2026,
\newblock \url{https://github.com/robopol/Riemann-hypothesis}.

\end{thebibliography}

\end{document}
